\chapter{特殊通道与专项计划}
\label{ch13}
\ist{特殊通道与专项}
\begin{keybox}
这一章讲推免主赛道之外的「岔路」。它们的共同特点是：名额常单列下达、条件写在部委文件里、一旦签约就带着服务期与违约条款。

按依据层级可分三类。\textbf{全国统一类}——研究生支教团（共青团中央办公厅、教育部办公厅2026年《意见》）、退役立功免初试与退役大学生士兵专项计划（《全国硕士研究生招生工作管理规定》第二十四、二十五条；教学厅〔2025〕8号）、少数民族骨干计划（硕士服务不少于5年）、公费师范生本研衔接（国办发〔2024〕27号）、优师计划（教师〔2021〕4号）、硕师计划（教师〔2009〕5号）、强基计划转段（教学〔2020〕1号）。\textbf{院校特定类}——强基转段的专业范围与是否占用招生计划、工程硕博士与国优计划的招生细则、各校对专升本与定向生的排除条款。\textbf{无统一规定类}——港澳台学生推免、分校区名额、中外合作办学海外成绩折算、定向选调与保研的冲突，本手册如实标注为「教育部层面未见统一规定，仅有各校做法」。

读这一章请记住三条：一是分清「专项计划」与「免初试」是两回事；二是凡涉及服务期、定向协议的，签约前逐条读清协议，文件规定只是底线；三是名额单列不等于门槛降低，多数专项对学业综合成绩的要求与普通推免一致。
\end{keybox}

\section{港澳台}
\begin{qa}{E0222}{退役后想保研，学校有没有单独的推免名额？}
结论：\textbf{教育部层面没有「退役生单独推免名额」这一制度，只有两条性质不同的路。}
\begin{itemize}
  \item 一是「退役大学生士兵专项硕士研究生招生计划」（教学厅〔2025〕8号）：这是\textbf{统考}招生计划，2026年由北京大学、清华大学等\textbf{532所}高校承担，在招生计划总规模内单列下达、专项使用，不得调整为普通计划。它\textbf{不给}推免资格。
  \item 二是各校推免细则里的口子：在校期间参军入伍、在部队荣立\textbf{二等功及以上}且已退役的应届本科毕业生，符合研究生报名条件的，学校可给予免试攻读研究生\textbf{推荐资格}（广东外语外贸大学2026年办法第十八条、湘潭大学2026年通知，后者援引教学厅〔2015〕3号）。这是\textbf{给资格}、占用本校推免名额，不是额外单列。
  \item 加分路径普遍存在但有限：多校把参军入伍服兵役纳入推免综合测评加分，但明文规定「\textbf{不得单列名额、单独排序}」，不得为仅符合单一指标的学生破格推荐（北京大学2027年办法、广东外语外贸大学2026年办法）。
\end{itemize}
\nref{S0134}
\nref{S0567}
\nref{S0052}
\nref{S0553}
\nref{S0079}
\end{qa}

\iin{广东外语外贸大学}
\ith{校内名额}
\ist{特殊通道与专项}
\begin{qa}{E0237}{国防生和军校生的推免政策是什么？跟普通生一样吗？}
结论：\textbf{不一样，要分三类看。}
\begin{itemize}
  \item \textbf{地方高校的国防生}：国防生制度已于2017年前后停止，在校国防生按政策分流；如你仍持「国防生」身份，请先向学校武装部与学院确认现行身份和协议义务，再谈升学。
  \item \textbf{军校生（生长军官学员）}：属军队院校，招生与研究生推免均走军队系统，不适用教育部面向地方高校的推免体系。地方高校推免办法的适用对象是「本校全日制应届本科毕业生」，覆盖不到军校学员。
  \item \textbf{地方高校应届生推免到军校}：这条路是通的。国防科技大学《接收2027级地方院校推荐免试硕士研究生（含直博生、参军入伍生）工作方案》要求申请人「具有推荐免试资格高校的应届本科毕业生，并取得所在学校硕士研究生推荐免试计划（\textbf{含我校『国防科工计划』}）」；另设「参军入伍生」类别，须符合军队直接选拔招录军官条件（本科院校为「双一流」建设高校、本科专业属指定学科门类、授予学位类型为法学/文学/理学/工学或管理学学士等），\textbf{具体政策以当年军队有关通知为准}。
\end{itemize}
\nref{S0650}
\nref{S0420}
\nref{S0572}
\nref{S0048}
\end{qa}

\iin{国防科技大学}
\ith{军校生}
\ist{特殊通道与专项}
\begin{qa}{E0246}{分校区（比如威海、深圳校区）的推免名额和本部一样吗？}
本所查公开资料中没有找到教育部关于「分校区推免名额如何分配」的统一规定，\textbf{教育部层面未见统一规定，仅有各校做法}。可以确认的是：名额下达的单位是\textbf{高校法人}，不是校区。
\begin{itemize}
  \item 制度层面：推免名额下达给「具有推免资格的高校」（《全国普通高等学校推荐优秀应届本科毕业生免试攻读硕士学位研究生工作管理办法（试行）》，教学〔2006〕14号），一个法人单位一个总数，校区之间怎么切由学校自己的名额分配办法决定。
  \item 校区在招生接收上是分列的：哈尔滨工业大学2027年接收推免生提醒中，校本部与\textbf{威海校区、深圳校区}以及未来工学院、地方研究院等并列作为接收单位并分别给出联系方式，说明接收环节各校区独立动作；但名额怎么分，该文件未给统一口径。
  \item 校区在制度文件上也分开：北京大学《关于推荐2027届优秀本科毕业生免试攻读研究生的办法》标题即注明「（\textbf{校本部}）」，说明校区/学部各有各的办法。
\end{itemize}
所以正确做法是：查本校《推免工作实施办法》的名额分配条款＋你所在校区/学院的当年细则，看具体指标数。
\nref{S0504}
\nref{S0553}
\nref{S0454}
\nref{S0568}
\end{qa}

\iin{哈尔滨工业大学}
\ith{分校区}
\ist{特殊通道与专项}
\begin{qa}{E0247}{分校区保研能不能保到本部或者别的校区？}
本所查公开资料中没有找到「分校区学生能否跨校区推免」的统一规定，\textbf{教育部层面未见统一规定，仅有各校做法}。制度上要把握三点：
\begin{itemize}
  \item \textbf{接收端是开放的}：推免录取是「招生单位—考生」双向选择，你在推免服务系统填志愿、对方发复试与待录取通知即可。哈尔滨工业大学2027年接收提醒把校本部与\textbf{威海、深圳校区}及未来工学院、地方研究院并列为接收单位，说明跨校区接收在操作上存在。
  \item \textbf{推荐端看本校名额办法}：推免资格从你本科所在学院/校区出，能否用该资格申请其他校区，取决于本校名额分配办法与学院是否放人；北京大学2027年推免办法标题即注明「（\textbf{校本部}）」，说明校区各有文件。
  \item \textbf{一条普遍约束}：多校推免办法明确「不得为仅符合单一或部分遴选指标的学生单列名额、单独排序」，同时要求各院系「积极推荐我校优秀毕业生到其他院校、科研机构攻读研究生，不囿于本校、本院系、本专业」。
\end{itemize}
所以先确认你手上有没有推免资格（这是前提），再按目标校区当年的接收细则投递。
\nref{S0504}
\nref{S0553}
\nref{S0454}
\end{qa}

\iin{哈尔滨工业大学}
\ith{分校区}
\ist{特殊通道与专项}
\begin{qa}{E0250}{第二学士学位在读能不能参加推免？}
结论：\textbf{第二学士学位学生不在推免对象范围内，全国高校口径高度一致。}
\begin{itemize}
  \item 制度源头：教育部《全国普通高等学校推荐优秀应届本科毕业生免试攻读硕士学位研究生工作管理办法（试行）》（教学〔2006〕14号）界定的对象是「普通高校应届本科毕业生」；2026年新印发的《全国高等学校推荐优秀应届本科毕业生免试攻读研究生工作管理规定》（\textbf{教学厅〔2026〕1号}）被各校援引后，推荐对象普遍表述为「本校通过普通高考（含保送方式）招生录取或面向港澳台华侨学生招生录取的全日制应届本科毕业生，\textbf{不含专升本、第二学士学位学生}」——成都大学商学院2027年细则、辽宁大学2027年办法、北京科技大学2027年通知、燕山大学2027年通知、甘肃农业大学2027年管理办法表述一致。
  \item 北京大学历年和2027年办法均写明「（\textbf{不含专升本、第二学士学位学生}）」；广东外语外贸大学2026年办法还额外排除辅修学位、本硕项目国际班学生和独立学院学生；湘潭大学2026年通知额外排除「职高对口」学生。
  \item 原因在于身份性质：第二学士学位属于「已获学士学位后的再读」，学籍按另一套办法管理，不属普通全日制应届本科毕业生。
\end{itemize}
第二学士学位想读研，请走统考。
\nref{S0454}
\nref{S0553}
\nref{S0554}
\nref{S0079}
\nref{S0168}
\nref{S0233}
\end{qa}

\iin{教育部学生服务与素质发展中心（学信网/研招网）}
\ith{第二学位}
\ist{特殊通道与专项}
\begin{qa}{E0252}{港澳台学生能不能参加内地高校的推免？}
结论：\textbf{分两种身份，答案不同。}
\begin{itemize}
  \item \textbf{在内地高校就读、通过面向港澳台华侨招生录取进来的学生}：\textbf{可以}参加推免。多校推荐对象明文写作「通过普通高考（含保送方式）招生录取或\textbf{面向港澳台华侨学生招生录取}的全日制本科生，不含专升本、第二学士学位学生」——成都大学商学院2027年细则、辽宁大学2027年办法、北京科技大学2027年通知、燕山大学2027年通知表述一致。
  \item \textbf{在港澳台本地院校或境外院校读本科的学生}：教育部层面\textbf{没有}统一的「港澳台学生推免办法」。推免的推荐环节只面向\textbf{本校}应届本科生，境外本科生没有推荐端；港澳台与外籍学生报考内地研究生走各校「港澳台招生/国际学生招生」渠道——中国人民大学研招网首页导航即把「硕士招生/推荐免试/港澳台生/国际学生」分设，各自有「简章目录、报名初试、复试录取」。
  \item 管理依据可佐证其独立性：《普通高等学校学生管理规定》（教育部令第41号）第六十六条规定「学校对接受高等学历继续教育的学生、\textbf{港澳台侨学生、留学生}的管理，参照本规定执行」。
\end{itemize}
\nref{S0233}
\nref{S0168}
\nref{S0079}
\nref{S0533}
\nref{S0285}
\nref{S0040}
\end{qa}

\iin{北京科技大学}
\ith{港澳台}
\ist{特殊通道与专项}
\begin{qa}{E0253}{港澳台生保研和内地生的政策有什么不同？}
结论：\textbf{教育部层面没有单独的「港澳台生推免政策」；差异出现在录取类别与报考渠道两个环节，而不是推免资格本身。}
\begin{itemize}
  \item \textbf{资格环节基本一致}：在内地高校就读、通过面向港澳台华侨招生录取进来的全日制应届本科生，与内地生适用同一套推荐条件与流程（成都大学商学院2027年细则、辽宁大学2027年办法、北京科技大学2027年通知、燕山大学2027年通知的推荐对象均把「面向港澳台华侨学生招生录取」写进范围）。
  \item \textbf{录取类别环节存在差异}：中山大学2027年推免招生录取工作办法写明「推免生的招生类别一般为\textbf{非定向就业}（录取为『少数民族高层次骨干人才计划』等国家专项计划的为定向就业）」，身份差别通常体现在定向/非定向与奖助安排上。
  \item \textbf{报考渠道环节差异明显}：境外/港澳台院校本科毕业生不走推免，而走各校「港澳台生」栏目（中国人民大学研招网单列「港澳台生-推荐免试/报名初试/复试录取」）。
  \item \textbf{管理依据另有层级}：教育部令第41号第六十六条明确港澳台侨学生、留学生的管理「参照本规定执行」。
\end{itemize}
\nref{S0167}
\nref{S0166}
\nref{S0233}
\nref{S0285}
\end{qa}

\iin{中山大学}
\ith{港澳台}
\ist{特殊通道与专项}
\begin{qa}{E0254}{港澳台生保研是不是要走专门的港澳台招生渠道？}
结论：\textbf{看你现在的学籍在哪，两种身份走两条完全不同的路。}
\begin{itemize}
  \item \textbf{已在内地高校读本科（通过普通高考或面向港澳台华侨招生录取）}：走\textbf{普通推免流程}，与内地同学一样在「全国推荐免试攻读研究生（免初试、转段）信息公开管理服务系统」注册、填报志愿、确认复试与待录取，\textbf{不需要}走「港澳台生」栏目。各校研招网的「港澳台招生」栏目是给另一类人用的：华南理工大学、上海交通大学、西南大学、中南大学等都在研招网单设「港澳台招生」入口。
  \item \textbf{在港澳台或境外院校读本科}：内地高校没有面向你的推免推荐渠道（推免的推荐端只面向本校应届本科生），你要走目标校的「\textbf{港澳台生招生}」栏目，按该校港澳台招生简章报名并参加考试。中国人民大学研招网即把「港澳台生」单列，下设「推荐免试/简章目录/报名初试/复试录取」，但那是该校面向港澳台学生的\textbf{独立招生办法}，不是推免服务系统。
  \item 走错渠道的直接后果是系统里报不上，所以先确认：你的本科是不是内地高校、是不是全日制普通应届。
\end{itemize}
\nref{S0481}
\nref{S0513}
\nref{S0632}
\nref{S0285}
\nref{S0468}
\end{qa}

\iin{华南理工大学}
\ith{港澳台}
\ist{特殊通道与专项}
\begin{qa}{E0260}{定向就业或者委培的本科毕业生能参加推免吗？}
结论：\textbf{定向就业/委培本科毕业生原则上不在推免对象内，除非原定向单位书面同意或协议已解除。}
\begin{itemize}
  \item 名额测算直接排除：安徽师范大学2027年推免方案规定，普通类基础指标按「全校2027届本科预计毕业学生（\textbf{不含专升本与定向就业生}）」总数折算。
  \item 推荐对象留有例外口子：浙江师范大学2027年推免通知的推荐对象为「全日制普通本科毕业生（\textbf{不含定向生，但原定向单位书面同意推荐或硕师计划有关教育局指定的除外}）」。
  \item 接收端也会查履约状态：中山大学2027年接收推免生报名通知与推免招生录取办法均写明「我校\textbf{不接收}未按国家有关规定完成履约责任（含违约）的『农村订单定向生』『教育部直属师范院校公费师范生』等推免生申请，一经查实，将取消其申请资格、录取资格、入学资格，已入学的取消学籍」。
  \item 服务期内的统考也受限：电子科技大学2026年章程规定，服务期限内的农村订单定向免费培养医学毕业生「\textbf{只能报考}临床医学专业学位硕士全科医学方向，且报考类别为\textbf{定向就业}」。
\end{itemize}
\textbf{签约前请逐条读清你的定向协议条款，以上为文件规定，个人协议可能另有约定。}
\nref{S0569}
\nref{S0646}
\nref{S0166}
\nref{S0167}
\nref{S0493}
\end{qa}

\iin{浙江师范大学}
\ith{定向委培}
\ist{特殊通道与专项}
\begin{qa}{E0263}{中医、中西医结合专业的推免有什么特殊通道？}
结论：\textbf{中医、中西医结合没有「专项推免通道」，但存在明确的专业限制条款，这是最容易踩的坑。}
\begin{itemize}
  \item 接收端的学科门槛（中山大学2027年推免招生录取工作办法）：「医学临床学科（1002临床医学、1003口腔医学、\textbf{1006中西医结合}）学术学位各专业\textbf{只接受授医学学位}的应届本科毕业生申请。」非医学学位的本科背景申请不了这些学术学位点。
  \item 学院层面各有细则：山东中医药大学中医学院发布《中医学院2027年优秀应届本科毕业生免试攻读研究生遴选工作实施细则》，其他中医药院校同样按学院发布细则，遴选条件、专业范围、名额都在学院文件里。
  \item 别把「5+3」当通道：西安交通大学2027年推免章程把「医学『5+3』一体化」列为专项计划，安徽医科大学2027年办法的对象为「推荐免试生及『5+3』一体化转段生」——那是长学制的转段生，不是给普通中医本科生开的通道。
\end{itemize}
建议：先看目标校医学类专业的「学位授予类型」要求，再对照本学院遴选细则，最后确认该专业是否接收外校推免生。
\nref{S0167}
\nref{S0639}
\nref{S0548}
\nref{S0600}
\end{qa}

\iin{中山大学}
\ith{医学专项}
\ist{特殊通道与专项}
\begin{qa}{E0266}{国家实验室联培、企业联培的推免项目是什么？}
结论：\textbf{联培类推免项目最主要的载体是「工程硕博士培养改革专项」，另有各校自设的实验室/企业联培项目；它们都以推免为入口、以校企（校院）双导师为特征。}
\begin{itemize}
  \item 组织方式（中国人民大学2027年工程硕博士培养改革专项推免研究生招生报名通知）：经批准在电子信息（085400）人工智能领域开展专项招生，分\textbf{工程硕士}与\textbf{工程直博}两个项目，实行「联合培养和工学交替模式，校企导师组共同承担研究生培养工作，学生深入工程攻关一线」；硕士基本学制3年（约1年校内课程＋2年联合培养单位实践），直博5年（约2年校内＋3年实践）；合作企业含腾讯、华为、百度、智谱华章、小米、北汽等。
  \item 招生机制：考生须有推免资格，报名走学校自己的系统（限报一个专业方向，不能同时申请硕士和直博项目），\textbf{最终录取仍须通过教育部推免服务系统}完成填报志愿、复试确认、待录取确认。
  \item 与普通推免的关系（东南大学2027年推免专项报名通知）：把「工程硕博士、集成电路专项班、国优计划、研支团」并列为四个专项，「\textbf{报名时只能选择一个专项，不可兼报}」「\textbf{确认拟录取的考生，不得再参选其他任何专项}」。
\end{itemize}
\nref{S0286}
\nref{S0186}
\nref{S0512}
\nref{S0188}
\nref{S0548}
\end{qa}

\iin{中国人民大学}
\ith{联培项目}
\ist{特殊通道与专项}
\begin{qa}{E0268}{艺术类、体育类的推免有专项通道吗？}
结论：\textbf{艺术类、体育类没有全国性的专项推免通道，通行做法是纳入普通推免的综合评价加分；少数学校设特长类推荐类型。}
\begin{itemize}
  \item 加分路径（湖南大学2027年公示）：《关于2027年推免生创新创业赛事（挑战杯）、体育类、艺术类竞赛获国家级及以上奖励、志愿服务加分名单的公示》按国家级名次折算分值——第22届全国大学生田径锦标赛4×400米混合接力第一名记8分、中国大学生跆拳道锦标赛（总决赛）品势混双第一名等合计11分、教育部全国大学生艺术展演（集体）国家级一等奖记1.67分。加分是「进入排名」，不是另给名额。
  \item 推荐类型路径：北京科技大学2027年通知把推荐类型分为普通类、\textbf{人才培养改革类}（理科试验班、工科试验班、黄昆英才班等）和研究生支教团等，并无艺体专项；东南大学2027年推免专项只列工程硕博士、集成电路专项班、国优计划、研支团四个。
  \item 上限约束：多校推免办法明确「\textbf{不得为仅符合单一或部分遴选指标的学生单列名额、单独排序}」（北京大学2027年办法、燕山大学2027年通知、广东外语外贸大学2026年办法均如此），纯靠竞赛加分撑起推免是有上限的。
\end{itemize}
\nref{S0541}
\nref{S0233}
\nref{S0512}
\nref{S0646}
\end{qa}

\iin{湖南大学}
\ith{艺术体育}
\ist{特殊通道与专项}
\begin{qa}{E0269}{定向选调生的报名和保研冲突了，怎么取舍？}
本所查公开资料中\textbf{没有检索到教育部关于「定向选调生报名与推免冲突」的任何规定性文件}。选调生是各省党委组织部门组织的地方公务员招录，不属于教育部研究生招生体系，两者之间没有全国统一的协调规则——\textbf{教育部层面未见统一规定，仅有各校做法}。
能确认的是推免一侧的硬约束：
\begin{itemize}
  \item 已被招生单位录取的推免生，\textbf{不得再报名参加当年全国统一考试、单独考试}（《2027年全国硕士研究生招生工作管理规定》第二十七条）；截止规定日期未被录取的推免生，其推免资格自动失效。
  \item 已落实接收单位的推免生，不得再申请出国留学、参加就业派遣等，否则违约责任自负（北京大学2027年推免办法）；已获得推免录取资格的申请人，应承诺\textbf{不再列入当年度就业计划}（西安交通大学2027年推免章程）。
  \item 时间上也容易撞车：2027年推免服务系统9月21日9:00开始填报志愿、10月18日截止填报志愿。
\end{itemize}
建议：把两条时间表（选调报名/笔试/政审 vs 推免填报与待录取确认）对齐，再逐字读本省选调公告里「应届毕业生」的界定，确认是否接受已获推免资格者。
\nref{S0442}
\nref{S0553}
\nref{S0548}
\end{qa}

\iin{教育部学生服务与素质发展中心（学信网/研招网）}
\ith{选调冲突}
\ist{特殊通道与专项}
\section{中外合作}
\begin{qa}{E0221}{退役免初试能不能和推免资格同时使用？}
结论：\textbf{两者互斥、不叠加；推免生本身就不用初试，不必也不应再去申请「免初试」。}
\begin{itemize}
  \item 「立功免初试」是考研报名环节政策（《2027年管理规定》第二十五条）：服役期间获三等战功、二等功以上奖励或二级以上表彰、符合报考条件的退役人员可申请免初试，须在统考报名期内登录「推免服务系统」报名，\textbf{逾期不得补报}，招生单位对审核通过者组织复试，复试合格即可拟录取。
  \item 硬约束：「\textbf{申请立功表彰免初试的考生，不得再报名参加全国硕士研究生招生考试}」（推免服务系统说明）；同时管理规定第二十七条明确「\textbf{已被招生单位录取的推免生，不得再报名参加当年全国统一考试、单独考试}」。
  \item 取舍很清楚：有推免资格就走推免（推免即免初试，还可填多个志愿）；没有推免资格但立过功，走「免初试」入口，代价是放弃当年统考。两者在同一个平台的不同入口，不能同时占。
  \item 另有连带约束：推免生一旦落实接收单位，不得再申请出国留学、参加就业派遣（北京大学2027年办法）。
\end{itemize}
\nref{S0442}
\nref{S0437}
\nref{S0464}
\nref{S0051}
\end{qa}

\iin{教育部学生服务与素质发展中心（学信网/研招网）}
\ith{政策叠加}
\ist{特殊通道与专项}
\begin{qa}{E0230}{地方专项计划会不会影响推免资格或者名额分配？}
结论：\textbf{专项计划（国家专项/地方专项/高校专项）是本科招生计划类型，不构成推免资格门槛；但部分学校会在名额测算时排除某些身份类别。}
\begin{itemize}
  \item 资格端：推免推荐对象的通行表述是「通过普通高考（含保送方式）招生录取或面向港澳台华侨学生招生录取的全日制应届本科毕业生，不含专升本、第二学士学位学生」（燕山大学2027年、北京科技大学2027年、成都大学商学院2027年表述一致），专项计划生属普通高考录取范畴，不被排除。
  \item 名额端：影响出现在「分母」和身份类别上。安徽师范大学2027年方案规定基础指标按「全校2027届本科预计毕业学生（\textbf{不含专升本与定向就业生}）」总数折算；浙江师范大学2027年通知的推荐对象「不含定向生，但原定向单位书面同意推荐或硕师计划有关教育局指定的除外」。
  \item 另一条独立通道：专项推免计划（研究生支教团、硕师计划等）名额由国家专项下达、单列分配，不参与按毕业生人数折算的普通名额（安徽师范大学2027年方案把「专项推免计划」与普通类并列）。
\end{itemize}
\nref{S0533}
\nref{S0233}
\nref{S0079}
\nref{S0646}
\nref{S0569}
\end{qa}

\iin{zsjyc（高校官网）}
\ith{地方专项}
\ist{特殊通道与专项}
\begin{qa}{E0244}{中外合作办学项目的学生，推免资格怎么算？}
结论：\textbf{资格本身不因「中外合作办学」被排除，但有两个变量决定你是否算在推免对象内：招生录取方式与学籍归属。}
\begin{itemize}
  \item 实践证据：浙江师范大学2027年免试攻读硕（博）士学位研究生\textbf{拟推荐名单公示}中，明确列有「小学教育（\textbf{中外合作办学}）」「音乐学（\textbf{中外合作办学}）」「软件工程（\textbf{中外合作办学}）」等专业的应届生，说明中外合作办学专业学生可在本校推免名额内被推荐。
  \item 排除条款针对的不是中外合作本身：各校推荐对象的排除项通常是「\textbf{不含专升本、第二学士学位学生}」（北京大学2027年办法、燕山大学2027年通知等）；广东外语外贸大学2026年办法排除范围更宽（辅修学位、专升本、第二学士学位、\textbf{本硕项目的国际班学生}和独立学院学生）。若你所在项目被学校归类为「国际班」或非普通全日制本科，才可能被排除。
  \item 机构层面已有先例：教育部办公厅2025年备案的67所新增推免资格高校中含\textbf{香港中文大学（深圳）、宁波诺丁汉大学}（首批中外合作办学），说明中外合作办学机构可整体获得推免资格。
\end{itemize}
先确认两件事：项目是否纳入国家普通本科招生计划的全日制本科；所在学校/校区本身有没有推免资格。
\nref{S0571}
\nref{S0640}
\nref{S0554}
\nref{S0134}
\end{qa}

\iin{jwc（高校官网）}
\ith{中外合作}
\ist{特殊通道与专项}
\begin{qa}{E0245}{中外合作办学在海外那两年的成绩算不算推免绩点？}
本所查公开资料中没有找到教育部关于「中外合作办学海外阶段成绩如何折算进推免绩点」的统一规定，\textbf{教育部层面未见统一规定，仅有各校做法}。可确认的相关事实：
\begin{itemize}
  \item 中外合作办学专业学生本身在推免名单里（浙江师范大学2027年免试攻读硕（博）士学位研究生拟推荐名单公示含多个「中外合作办学」专业），说明其课程成绩必然被纳入某种计算体系。
  \item 推免遴选以「学业综合成绩」为基础指标，成绩口径由学校教务部门规定。北京大学2027年办法写明「学生成绩数据以\textbf{教务系统统计为准}」「\textbf{各院系的综合排名标准由院系自行制定和解释}」，说明不进教务系统的成绩通常不进排名。
  \item 交换/双学位成绩通常走学校\textbf{学分认定与成绩转换}流程（教务处学分认定办法），转换规则、是否计入绩点、按几分制折算各校不同；有的学校只认学分不认绩点。
  \item 一条普遍约束：多校明文规定「不得为仅符合单一或部分遴选指标的学生单列名额、单独排序」，因此不存在「海外成绩单独加分」。
\end{itemize}
正确做法：向教务处索取本校《本科生课程学分认定与成绩转换办法》。
\nref{S0571}
\nref{S0554}
\nref{S0633}
\end{qa}

\iin{jwc（高校官网）}
\ith{中外合作}
\ist{特殊通道与专项}
\begin{qa}{E0248}{中外合作办学的学生能不能保研到国内其他学校？}
结论：\textbf{只要你有推免资格，保研到哪所学校与「中外合作办学」身份没有制度性关系；是否受限取决于本校有无附加限制与目标校报考条件。}
\begin{itemize}
  \item 制度导向（北京大学2027年推免办法）：各院系要「积极推荐我校优秀毕业生到\textbf{其他院校、科研机构}攻读研究生，不囿于本校、本院系、本专业，鼓励学生积极联系外单位」。这是全国普遍的工作导向，与本科专业是否中外合作无关。
  \item 实操上落到你头上的是两件事：一是本校/本院细则里有没有对中外合作办学专业设「仅限本校」「仅限项目内」的限制（属院校特定，需查本院细则）；二是目标院校/专业对本科背景的报考条件是否设限。
  \item 已有实践证据：浙江师范大学2027年拟推荐名单中中外合作办学专业学生被列为推免生，说明他们在同一套推荐体系内、可外推。
  \item 需注意：部分中外合作项目在协议中对出国阶段或学位授予有附加要求，若与推免时间表、去向绑定，会影响你在10月18日填报志愿截止前的选择。
\end{itemize}
\nref{S0553}
\nref{S0571}
\nref{S0633}
\nref{S0554}
\end{qa}

\iin{北京大学}
\ith{中外合作}
\ist{特殊通道与专项}
\begin{qa}{E0249}{专升本的学生到底有没有推免资格？}
结论：\textbf{没有。全国高校口径高度一致：专升本不在推免推荐对象范围内。}
\begin{itemize}
  \item 学校原文几乎是同一个模板：「申请推免资格的学生应为应届本科毕业生，系本校通过普通高考（含保送方式）招生录取或面向港澳台华侨学生招生录取的全日制本科生，\textbf{不含专升本、第二学士学位学生}」——燕山大学2027年通知、北京科技大学2027年通知、辽宁大学2027年办法、成都大学商学院2027年细则、甘肃农业大学2027年管理办法均如此表述。
  \item 北京大学2027年办法与2017年实施办法亦写明：申请推免学生须「纳入国家普通本科招生计划录取的应届毕业生（\textbf{不含专升本、第二学士学位学生}）」。
  \item 更宽的排除还有：广东外语外贸大学2026年办法排除辅修学位、专升本、第二学士学位学生、本硕项目国际班学生和独立学院学生；湘潭大学2026年通知额外排除「职高对口」学生。
  \item 医学类院校再排一层：温州医科大学2026、2027年通知的推荐对象「不含专升本、第二学士学位、\textbf{5+3一体化、订单定向培养学生}」。
\end{itemize}
所以专升本想读研请走统考路线，不要把时间花在推免准备上。
\nref{S0533}
\nref{S0233}
\nref{S0553}
\nref{S0554}
\nref{S0134}
\nref{S0567}
\end{qa}

\iin{zsjyc（高校官网）}
\ith{专升本}
\ist{特殊通道与专项}
\begin{qa}{E0251}{专升本加考研的成绩怎么算？推免认可专科阶段的成绩吗？}
结论：\textbf{推免端不认可专升本，因为专升本学生整体不在推免推荐对象内；统考端不受此限，本科阶段成绩照常使用。}
\begin{itemize}
  \item 推免端：各校推荐对象都是「本校通过普通高考（含保送方式）招生录取或面向港澳台华侨学生招生录取的\textbf{全日制应届本科生}，\textbf{不含专升本、第二学士学位学生}」（燕山大学2027年、北京科技大学2027年、辽宁大学2027年表述一致）——门槛在\textbf{身份}上就已排除，不存在「成绩怎么算」的问题，专科阶段成绩更不会被纳入。
  \item 统考端不歧视：福州大学2026年招生简章按国家规定把「国家承认学历的应届本科毕业生（含普通高校、成人高校、普通高校举办的成人高等学历教育等应届本科毕业生）及自学考试和网络教育届时可毕业本科生」列为可报考对象，条件是「录取当年入学前必须取得国家承认的本科毕业证书或教育部留学服务中心出具的《国（境）外学历学位认证书》，否则录取资格无效」（福建农林大学2026年招生简章为同一表述）。
\end{itemize}
所以建议把精力放在统考初试与复试，把本科两年成绩单、四六级、科研经历准备好。
\nref{S0533}
\nref{S0233}
\nref{S0603}
\nref{S0413}
\end{qa}

\iin{zsjyc（高校官网）}
\ith{专升本}
\ist{特殊通道与专项}
\begin{qa}{E0255}{国际学生（留学生）能不能保研到中国的高校？}
结论：\textbf{外籍国际学生不走推免，走各校「国际学生」招生渠道；该渠道在教育部层面没有统一的「推免」规定。}
\begin{itemize}
  \item 推荐端已排除：推免推荐对象是「本校全日制应届本科毕业生」，多校进一步限定为「通过普通高考（含保送方式）招生录取或面向港澳台华侨学生招生录取」的全日制本科生（燕山大学2027年、北京科技大学2027年、辽宁大学2027年表述一致）——外籍学生不在此列。
  \item 接收端分设栏目：各校研究生院把「港澳台生」「国际学生」单列为独立栏目，如中国人民大学研招网首页导航分设「硕士招生/推荐免试/港澳台生/国际学生」，各自设「简章目录、报名初试、复试录取」，说明是按单独的招生办法组织考试录取，不是推免。
  \item 学籍管理依据：《普通高等学校学生管理规定》（教育部令第41号）第六十六条规定「学校对接受高等学历继续教育的学生、\textbf{港澳台侨学生、留学生}的管理，参照本规定执行」，制度上把留学生单独归类。
\end{itemize}
所以外籍学生想读中国高校研究生，请查目标校「国际学生招生简章」，按申请审核或考试方式准备。
\nref{S0285}
\nref{S0533}
\nref{S0233}
\nref{S0040}
\nref{S0366}
\end{qa}

\iin{中国人民大学}
\ith{留学生}
\ist{特殊通道与专项}
\begin{qa}{E0256}{交换生、双学位学生的海外成绩怎么折算进推免排名？}
本所查公开资料中没有教育部关于「交换生/双学位海外成绩如何折算进推免排名」的统一规定，\textbf{教育部层面未见统一规定，仅有各校做法}。可把握的框架是：
\begin{itemize}
  \item 推免遴选以「学业综合成绩」为基础指标，成绩口径由学校教务系统决定。北京大学2027年办法写明「学生成绩数据\textbf{以教务系统统计为准}」「\textbf{各院系的综合排名标准由院系自行制定和解释}」——海外成绩进不进排名，取决于它有没有被教务系统认定并录入。
  \item 交换/双学位成绩通常走学校的\textbf{学分认定与成绩转换}流程（教务处学分认定办法），转换规则、是否计入绩点、按几分制折算各校不同；有的学校只认学分、不认绩点。
  \item 一条普遍约束：多校明文规定「不得为仅符合单一或部分遴选指标的学生单列名额、单独排序」，因此不存在「海外成绩可以单列加分」的做法。
  \item 成绩结构会影响结果：燕山大学、内蒙古大学等把综合成绩写成「加权平均学分绩点＋竞赛加分＋参军入伍加分」，海外成绩若无法折算成绩点，就无法进入这一公式。
\end{itemize}
\nref{S0554}
\nref{S0553}
\nref{S0544}
\nref{S0533}
\end{qa}

\iin{北京大学}
\ith{留学交换}
\ist{特殊通道与专项}
\begin{qa}{E0257}{支教团、少干、退役这几类专项能不能同时申请？}
结论：\textbf{推免类专项之间基本不能兼报；退役「立功免初试」不属于推免专项，另行报名且与统考互斥。}
\begin{itemize}
  \item 系统层级（研招网「考试方式、专项计划概述以及相互关系」）：「当考试方式选择\textbf{推荐免试}时，考生可选择专项计划为『无』、『\textbf{少数民族骨干计划}』、『\textbf{支教团推免计划}』或者『\textbf{免费师范生计划}』」——这些是推免方式下的并列选项，\textbf{一次只能选一个}。
  \item 学校层级（东南大学2027年推免专项报名通知）：「报名时\textbf{只能选择一个专项，不可兼报}……提交后不能再做更改」；「\textbf{确认拟录取的考生，不得再参选其他任何专项}」。
  \item 支教团会占用你的推免身份：安徽工业大学2027年支教团方案要求报名时同时提交《推荐免试攻读研究生资格申请表》，并在学生推免系统完成信息录入。
  \item 少干的互斥更硬：骨干计划考生「不得兼报其他专项计划，不得兼报硕士研究生普通计划」（贵州省2026年通知）。
  \item 退役路径另算：立功免初试在推免服务系统的「免初试」入口报名，且「\textbf{申请立功表彰免初试的考生，不得再报名参加全国硕士研究生招生考试}」。
\end{itemize}
\nref{S0449}
\nref{S0512}
\nref{S0589}
\nref{S0573}
\nref{S0464}
\end{qa}

\iin{教育部学生服务与素质发展中心（学信网/研招网）}
\ith{政策叠加}
\ist{特殊通道与专项}
\begin{qa}{E0261}{农村订单定向医学生能保研吗？服务期怎么处理？}
结论：\textbf{原则上不能——这是接收端有明文排斥条款的少数情形之一；例外是已完成履约责任或取得定向单位与卫健部门书面同意。}
\begin{itemize}
  \item 接收端直接拒绝（中山大学2027年接收推免生报名通知与推免招生录取工作办法）：「我校\textbf{不接收}未按国家有关规定完成履约责任（含违约）的『\textbf{农村订单定向生}』『教育部直属师范院校公费师范生』等推免生申请，一经查实，将取消其申请资格、录取资格、入学资格，已入学的取消学籍。」
  \item 推荐端也常直接排除：温州医科大学2026、2027年推荐资格认定通知的对象为「纳入国家普通本科招生计划录取的应届毕业生（\textbf{不含专升本、第二学士学位、5+3一体化、订单定向培养学生}）」。
  \item 即便走统考，服务期内也受限：电子科技大学2026年招生章程规定「服务期限内的农村订单定向免费培养医学毕业生\textbf{只能报考临床医学专业学位硕士全科医学方向}，且报考类别为\textbf{定向就业}；考生在报考时必须提供定向就业县（市、区）卫生健康行政部门出具的同意报考\textbf{书面同意函}」。
\end{itemize}
\textbf{签约前请逐条读清你的定向协议条款，以上为文件规定，个人协议可能另有约定。}
\nref{S0166}
\nref{S0167}
\nref{S0562}
\nref{S0563}
\nref{S0493}
\end{qa}

\iin{中山大学}
\ith{订单定向}
\ist{特殊通道与专项}
\begin{qa}{E0262}{医学的 5+3 一体化算不算推免？跟保研有什么区别？}
结论：\textbf{不算推免，是「转段」，两者是并行的不同机制。}
\begin{itemize}
  \item 制度定性（西安交通大学2027年推荐免试研究生招生章程）：把「\textbf{医学『5+3』一体化}」与强基计划、国优计划、工程硕博士专项、支教团、退役士兵专项并列列为「\textbf{专项计划}」，并写明「根据上级部门要求执行，学校根据工作安排另行通知」——学校按专项办，但走的是转段流程。
  \item 招生方式上两类人分列：安徽医科大学2027年办法的对象写作「推荐免试生\textbf{及『5+3』一体化转段生}」，并分别发布《接收推免生及『5+3』一体化转段生专业目录》《推荐免试生及『5+3』一体化转段生政治审查表》。
  \item 身份差异：推免是面向外校的招生选拔，需本校推免资格并在推免服务系统填报志愿；5+3转段是本校长学制培养方案内的阶段转换，通常不占用对外推免名额、也不需要另找接收单位。
  \item 独立性可由排除条款佐证：温州医科大学2026、2027年推免推荐资格认定通知把「\textbf{5+3一体化}」与专升本、第二学士学位、订单定向生一起排除在推荐对象之外——正因为它不属于推免体系。
\end{itemize}
\textbf{涉及服务期与协议条款的，签约前请逐条读清协议，以上为文件规定。}
\nref{S0548}
\nref{S0600}
\nref{S0563}
\nref{S0562}
\end{qa}

\iin{西安交通大学}
\ith{医学学制}
\ist{特殊通道与专项}
\section{强基转段}
\begin{qa}{E0223}{强基计划转段怎么确认？10 月 13 到 20 号这段时间要做什么？}
结论：\textbf{10月13日—20日是「推免服务系统」里强基转段信息的查询确认窗口，这段时间只做一件事：核对并确认本人转段信息。}
\begin{itemize}
  \item 官方原文（「全国推荐免试攻读研究生（免初试、转段）信息公开管理服务系统」说明第（六）条）：「\textbf{2023级拟转段的强基生于2026年10月13日-10月20日通过『推免服务系统』查询并确认本人转段信息，确认后的信息将作为强基生进入研究生阶段的录取信息。}」
  \item 同期时间轴（同页）：2027年推免生自2026年9月18日起注册填报资料、查看本人推免资格；9月21日—10月20日填报志愿、确认复试及待录取通知（9月21日9:00开始填报志愿，\textbf{10月18日截止填报志愿}；9月24日14:00开始确认待录取通知）。
  \item 校内流程可作对照（华中科技大学2026年强基转段通知）：7月13日前在研究生招生管理系统提交攻读意向（\textbf{每人限选1个专业}）→7月14—16日意向院系面试考核→院系接收结果7月16日前上报研究生院→研究生院与本科生院汇总核对后上报教育部研招网。
\end{itemize}
\nref{S0464}
\nref{S0536}
\nref{S0248}
\end{qa}

\iin{教育部学生服务与素质发展中心（学信网/研招网）}
\ith{强基转段}
\ist{特殊通道与专项}
\begin{qa}{E0224}{强基转段不想读原来的专业了，能转出去吗？}
结论：\textbf{可以在强基计划规定的专业范围内跨专业甚至跨学科申请，但不能转到范围之外的任意专业，也不能转出校。}
\begin{itemize}
  \item 华中科技大学2026年强基转段通知：「获得转段资格的强基生，\textbf{只能申请校内相关院系}的硕士或直接攻读博士」；「可根据所在专业强基计划培养方案和招生专业范围，申请在本学科相关专业或\textbf{跨学科}继续深造。跨学科交叉培养的\textbf{专业范围应符合学校和所在院系公布的要求}」；报名时每人限选1个专业，跨院系申请未通过考核的，须参加本科院系考核确定最终专业意向。
  \item 专业范围有清单：中山大学生命科学学院2026年转段方案列出0710生物学、0713生态学、0860生物与医药、1001基础医学、0812计算机科学与技术、0854电子信息等，本年度按\textbf{直博生}录取，「博士研究生指标将\textbf{单列下达}」。
  \item 制度底色（教学〔2020〕1号）：强基计划「探索建立本—硕—博衔接的培养模式，硕博阶段既可在本学科深造，也可探索\textbf{学科交叉培养}」。
\end{itemize}
\nref{S0536}
\nref{S0248}
\nref{S0045}
\end{qa}

\iin{华中科技大学}
\ith{强基转段}
\ist{特殊通道与专项}
\begin{qa}{E0225}{强基转段能不能放弃，去参加普通推免或者考研？}
结论：\textbf{可以放弃转段，但代价是失去推免资格；且获得转段拟录取资格后，出国与就业推荐服务会被暂停。}
\begin{itemize}
  \item 放弃的直接后果（中山大学生命科学学院2026年转段方案）：「\textbf{未参加或未通过转段考核的强基计划本科毕业生，视为退出强基计划，不再具有申请免试攻读研究生资格}」；「在本研转段期间退出强基计划的学生，由本院（系）按照\textbf{普通班培养方案}组织学分认定」。即：退出=同时失去推免。
  \item 锁定条款（华中科技大学2026年通知）：「学生报名申请转段后，一经审核确定获得转段生拟录取资格的，\textbf{原则上不得再办理出国和毕业派遣手续}」；「如2027年未能按期毕业，视作自动失去拟录取资格」。中山大学方案亦写明获得转段拟录取资格者，学校将暂停提供包括办理成绩单在内的出国留学证明及本科毕业就业推荐等服务。
  \item 考研：教育部层面未禁止退出强基的学生统考，但你要按普通班重新认定学分，且失去推免资格后只能统考。
\end{itemize}
\nref{S0248}
\nref{S0536}
\nref{S0045}
\end{qa}

\iin{中山大学}
\ith{强基转段}
\ist{特殊通道与专项}
\begin{qa}{E0226}{强基转段没通过怎么办？还有别的出路吗？}
结论：\textbf{转段考核没通过=退出强基计划、不再具有推免资格；出路只剩统考、出国或就业，且出国与就业推荐服务在转段期间会被暂停。}
\begin{itemize}
  \item 直接后果（中山大学生命科学学院2026年转段方案）：「未参加或未通过转段考核的强基计划本科毕业生，视为\textbf{退出强基计划，不再具有申请免试攻读研究生资格}」；同时「由本院（系）按照普通班培养方案组织\textbf{学分认定}」。
  \item 会「不通过」的具体情形（可提前规避）：华中科技大学2026年通知规定复试成绩不及格（百分制60分以下）不予接收、思想政治素质和品德考核不合格不予接收；中山大学方案规定复试总成绩或各分项任一项低于60\%为不合格，另有体检不合格、违纪受处分、\textbf{本科最后一学年课程成绩不及格}、不能按期毕业或未获学士学位等取消拟录取资格的情形。
  \item 所以真实出路是提前守住可控项：最后一学年不出现不及格、按期毕业拿学位、避免违纪，同时把统考作为保底方案并行准备。
\end{itemize}
\nref{S0248}
\nref{S0536}
\nref{S0045}
\end{qa}

\iin{中山大学}
\ith{强基转段}
\ist{特殊通道与专项}
\begin{qa}{E0227}{强基转段是直博还是硕士？能自己选吗？}
结论：\textbf{硕士和直博两种都可能，但不少学校近年按「直博生」录取；层次由院系当年方案与指标决定，不是单方面可选。}
\begin{itemize}
  \item 以直博为主：中山大学生命科学学院2026年转段方案明确「本年度强基生转段\textbf{按照直博生录取}，博士研究生指标将\textbf{单列下达}，最终招生人数以实际招收人数为准」。
  \item 两种都可选：华中科技大学2026年通知「获得转段资格的强基生，只能申请校内相关院系的\textbf{硕士或直接攻读博士}」，报名时提交攻读意向（含意向学院、专业、\textbf{攻读学位层次}），每人限选1个专业；跨院系申请未通过考核的，须参加本科院系考核确定最终专业意向。
  \item 政策给的是方向而非定式（教学〔2020〕1号）：强基计划「探索建立\textbf{本—硕—博衔接}的培养模式……高校可在免试推荐研究生、\textbf{直博}、公派留学、奖学金等方面予以优先安排」，具体层次由各校培养方案决定。
\end{itemize}
\nref{S0248}
\nref{S0536}
\nref{S0045}
\nref{S0368}
\end{qa}

\iin{中山大学}
\ith{强基转段}
\ist{特殊通道与专项}
\begin{qa}{E0242}{本硕博贯通、拔尖计划实验班怎么进去？进去了还需要自己保研吗？}
结论：\textbf{进入方式以高考录取和入校后的校内选拔为主；进班不等于自动读研，多数情况仍要走一次考核或转段确认。}
\begin{itemize}
  \item 拔尖班仍要申请推免：北京科技大学2027年推免通知把「人才培养改革类（理科试验班、工科试验班、黄昆英才班、钢铁脊梁计划、科史新星计划、青马人才工程、材料高精尖班、机器人科创班、卓越创新班、纳米试验班）」与普通类、研究生支教团并列为\textbf{推荐类型}，并明确「不同类型有不同的推荐条件……\textbf{学生申请其中一项推荐类型}」。
  \item 贯通不等于免确认：强基转段生仍须在推免服务系统确认信息（2023级拟转段强基生于2026年10月13日—20日查询并确认本人转段信息）；学校层面还要走意向提交、院系考核、公示上报。
  \item 强基计划的制度逻辑（教学〔2020〕1号）：「探索建立\textbf{本—硕—博衔接}的培养模式……建立科学化、多阶段的\textbf{动态进出机制}，对进入强基计划的学生进行综合考查、科学分流」，即贯通培养伴随分流淘汰。
\end{itemize}
\nref{S0233}
\nref{S0045}
\nref{S0548}
\nref{S0464}
\end{qa}

\iin{北京科技大学}
\ith{拔尖计划}
\ist{特殊通道与专项}
\begin{qa}{E0243}{强基、拔尖班的学生能不能放弃转段去参加普通推免？}
结论：\textbf{强基生不能「既放弃转段又走普通推免」——放弃转段的直接后果就是失去申请免试攻读研究生的资格。}
\begin{itemize}
  \item 明确原文（中山大学生命科学学院2026年强基计划本研转段工作实施方案）：「\textbf{未参加或未通过转段考核的强基计划本科毕业生，视为退出强基计划，不再具有申请免试攻读研究生资格}」；「在本研转段期间退出强基计划的学生，由本院（系）按照普通班培养方案组织学分认定」。
  \item 有约束的时点（华中科技大学2026年强基转段通知）：「学生报名申请转段后，一经审核确定获得转段生拟录取资格的，\textbf{原则上不得再办理出国和毕业派遣手续}」；中山大学方案亦写明获得转段拟录取资格者，学校将暂停提供包括办理成绩单在内的出国留学证明及本科毕业就业推荐等服务。
  \item 对\textbf{非强基的拔尖班}：通常只是推荐类型不同，普通推免通道仍然有效，具体看本校推免办法里各类推荐类型的申请规则（北京科技大学2027年通知即区分普通类、人才培养改革类、研究生支教团，学生申请其中一项推荐类型）。
\end{itemize}
\nref{S0248}
\nref{S0536}
\nref{S0233}
\nref{S0045}
\end{qa}

\iin{中山大学}
\ith{拔尖计划}
\ist{特殊通道与专项}
\begin{qa}{E0258}{强基转段信息确认之后还能改吗？}
结论：\textbf{系统里「确认」这一步基本不可逆——确认后的信息即作为录取信息；要调整只能在此前的校内环节改。}
\begin{itemize}
  \item 系统口径（推免服务系统说明第（六）条）：「2023级拟转段的强基生于2026年10月13日-10月20日通过『推免服务系统』查询并\textbf{确认}本人转段信息，\textbf{确认后的信息将作为强基生进入研究生阶段的录取信息}。」其实质是把确认结果直接固化为录取数据。
  \item 可调整的空间在前面（华中科技大学2026年强基转段通知流程）：在研究生招生管理系统提交攻读意向（学院、专业、攻读学位层次，\textbf{每人限选1个专业}）→意向院系考核面试（7月14—16日）→若跨院系申请未通过考核，\textbf{须参加本科院系的考核确定最终专业意向}→各接收院系接收结果7月16日前上报研究生院→研究生院与本科生院汇总核对后上报教育部研招网。可见改动只发生在7月的意向提交与考核阶段。
  \item 另有一类「确认后仍会被取消拟录取资格」的情形：信息不真实或伪造材料、未按规定报备回避关系、复试不及格、体检不合格、最后一学年课程不及格、不能按期毕业或未获学士学位等——这时不是「改」，而是「作废」。
\end{itemize}
\nref{S0464}
\nref{S0536}
\nref{S0248}
\end{qa}

\iin{教育部学生服务与素质发展中心（学信网/研招网）}
\ith{强基转段}
\ist{特殊通道与专项}
\section{少干计划}
\begin{qa}{E0228}{硕师计划是什么？怎么申请？要服务几年？}
结论：\textbf{硕师计划=从有推免资格的高校选拔应届本科生免试攻读教育硕士，同时被聘为编制内教师，到县镇及以下农村学校任教3年、第4年脱产学习。}
\begin{itemize}
  \item 政策原文（教育部《关于做好2010年「农村学校教育硕士师资培养计划」实施工作的通知》，教师〔2009〕5号）：「从具有推荐免试硕士研究生资格的高校中，选拔部分优秀应届普通本科毕业生，录取为『硕师计划』研究生，并与地方政府教育行政部门\textbf{签约聘为编制内正式教师}。在\textbf{县镇及以下农村学校任教，服务期三年}，并在职学习研究生课程。\textbf{第四年，到培养学校脱产集中学习一年}，毕业时获硕士研究生毕业证书和教育硕士专业学位证书。」
  \item 申请流程：各校发布年度实施通知→学生报名→学院资格审查→与县（市、区）教育行政部门及用人学校\textbf{供需见面/双选会}→签约→复试。吉林师范大学2027年通知要求「志愿到县镇及以下农村学校任教」、毕业时取得本科毕业证、学位证和教师资格证。
  \item 谁不能报：吉林师范大学2027年通知推荐对象「不含专升本、第二学士学位及免费师范生」；云南师范大学2025年办法「不含专升本、第二学士学位的学生、免费定向师范生」。
  \item 新变化：2027年起部分省份把「优师计划」与「硕师计划」结合实施。
\end{itemize}
\nref{S0525}
\nref{S0593}
\nref{S0377}
\nref{S0171}
\nref{S0592}
\end{qa}

\iin{教育部学生服务与素质发展中心（学信网/研招网）}
\ith{硕师计划}
\ist{特殊通道与专项}
\begin{qa}{E0231}{少数民族骨干计划怎么报？需要准备哪些材料？}
结论：\textbf{分两步走：先在「骨干计划数字化管理平台」做资格审核，通过后才能上研招网报名；没通过审核就报不上。}
\begin{itemize}
  \item 流程（贵州省教育厅2026年通知，依据教民厅〔2025〕1号）：须通过研招网「少数民族高层次骨干人才研究生招生计划资格申请平台」（http://mz.chsi.com.cn/mzjh/stu/）提交申请，经\textbf{生源地省级教育行政部门在线审核通过后方可报名}；平台开放时间为2025年10月9日9:00至10月24日22:00，工作人员3个工作日内完成审核。
  \item 所需材料（贵州原文）：1.身份证国徽面和人像面；2.居民户口簿首页和本人页；3.学历证书（往届）或就读证明/教育部学籍在线验证报告（应届）；4.学位证（往届）；5.工作单位同意报考证明（在职必提）；6.在职证明材料（在职必提，汉族考生须为本地户籍且出具在民族自治地方3年以上在职证明）。格式支持jpg/jpeg/png/pdf，每类上限3个，总大小10M以内。
  \item 签约与规模：被录取者通过骨干计划数字化管理平台\textbf{在线签署定向就业协议书}；2026年教育部安排博士招生计划1500名（名额不分配到省）、硕士5000名，其中贵州338名。
\end{itemize}
\nref{S0573}
\nref{S0595}
\nref{S0578}
\end{qa}

\iin{中国政府网}
\ith{少干计划}
\ist{特殊通道与专项}
\begin{qa}{E0232}{少干计划是走推免还是统考？能不能推免走少干？}
结论：\textbf{两条都能走，推免可以走少干。}
\begin{itemize}
  \item 系统层级（研招网「考试方式、专项计划概述以及相互关系」）：「当考试方式选择\textbf{推荐免试}时，考生可选择专项计划为『无』、『\textbf{少数民族骨干计划}』、『支教团推免计划』或者『免费师范生计划』。」说明少干是推免方式下的可选专项。
  \item 学校层级：中国人民大学2025年骨干计划招生简章中，硕士报考条件明确列有「符合条件的\textbf{推荐免试生可申请本计划}」，同时要求「承诺毕业后履行定向协议回定向地区和单位就业」，且不接受同等学力考生。
  \item 统考路径同样成立（贵州省2026年通知）：骨干计划硕士考生参加全国硕士研究生招生考试，实行「\textbf{自愿报考、统一考试、单独划线、择优录取}」，由教育部按学科门类分别确定进入复试的初试成绩基本要求。
  \item 互斥条款：骨干计划考生「不再享受其他初试加分和照顾政策，\textbf{不得兼报其他专项计划，不得兼报硕士研究生普通计划}」；自2028年起，原少数民族照顾政策与骨干计划整合为「少数民族和民族地区高层次骨干人才研究生招生计划」（2027年管理规定第二十三条）。
\end{itemize}
\nref{S0449}
\nref{S0578}
\nref{S0573}
\nref{S0442}
\end{qa}

\iin{教育部学生服务与素质发展中心（学信网/研招网）}
\ith{少干计划}
\ist{特殊通道与专项}
\begin{qa}{E0233}{少干计划毕业后必须回原籍工作吗？要服务几年？}
结论：\textbf{必须履行定向协议就业。硕士服务年限不得少于5年（西藏班、新疆班教师和管理人员为8年），博士不少于8年；文件口径是「考生不得违约」。}
\begin{itemize}
  \item 原文（贵州省教育厅2026年通知，依据教民厅〔2025〕1号）：「骨干计划\textbf{在职考生回本人原工作单位就业}；非在职考生\textbf{必须在定向省份或内蒙古、广西、贵州、云南、西藏、青海、宁夏、新疆（含兵团）等定向地区就业，不包分配工作}。硕士毕业服务年限不得少于5年（含5年，其中\textbf{西藏班、新疆班教师和管理人员为8年}），博士毕业服务年限不得少于8年（含8年）。\textbf{考生不得违约}。」
  \item 在学期间的连带限制：未经生源地省级教育行政部门和在职研究生原工作单位同意，骨干计划硕士在学期间和服务期内\textbf{不得报考博士研究生}；经同意可报考骨干计划博士并签新协议，毕业后服务年限\textbf{按新协议重新计算}。
  \item 证书与档案流向：非在职考生的毕业证、学位证和学生档案由招生单位寄送至省级毕业生就业工作部门，不在本人手上。
\end{itemize}
\textbf{签约前请逐条读清协议条款，以上为文件规定，个人协议可能另有约定。}
\nref{S0573}
\nref{S0595}
\nref{S0578}
\end{qa}

\iin{中国政府网}
\ith{少干计划}
\ist{特殊通道与专项}
\begin{qa}{E0234}{少干定向协议怎么签？违约了会怎么样？}
结论：\textbf{协议通过「骨干计划数字化管理平台」在线签署，是三方（或四方）定向就业协议；文件口径是「考生不得违约」，但违约金标准在协议里约定，全国没有统一依据。}
\begin{itemize}
  \item 签署方式（贵州省教育厅2026年通知）：被录取考生通过\textbf{骨干计划数字化管理平台}签署《少数民族高层次骨干人才计划毕业生定向就业协议书》，「签订协议时间待教育部与招生单位通知」；山东省说明为「在线签订\textbf{三方（或四方）定向就业协议书}，毕业后履行定向协议回定向地区和单位就业」。
  \item 签约前的资信审查：贵州省通知明确硕士资格审查通过名单须在平台\textbf{公示}，「凡弄虚作假者，一经查实，一律取消报考和录取资格」。
  \item 违约后果的边界：教育部层面明确「\textbf{考生不得违约}」、服务年限不得少于5年（博士8年、西藏班新疆班教师8年），但\textbf{具体违约责任由协议书约定}；本所查公开资料\textbf{未检索到}全国统一的骨干计划违约金计算标准。
  \item 可参照的同类机制：公费师范生违约处理为「按协议约定承担相应责任，违约行为\textbf{纳入信用记录}，并与全国信用信息共享平台共享，违约相关材料归入本人人事档案」（国办发〔2024〕27号第十三条）。
\end{itemize}
\textbf{签约前请逐条读清协议条款，以下为文件规定，个人协议可能另有约定。}
\nref{S0573}
\nref{S0595}
\nref{S0578}
\nref{S0456}
\end{qa}

\iin{中国政府网}
\ith{少干计划}
\ist{特殊通道与专项}
\begin{qa}{E0235}{国防军工定向生怎么保研？需要政审吗？}
结论：\textbf{「国防科工」类推免是学校推免计划中的一个定向补偿类型，接收端主要是军队院校与国防单位；思想政治素质和品德考核（政审）是必备环节。}
\begin{itemize}
  \item 计划性质：南开大学推免实施办法把「\textbf{定向补偿推免（支教团、硕师计划、国防科工院校补偿计划、新疆高层次双语人才计划等）}」单列，并规定「工作由牵头部门负责组织实施」——属专项补偿计划，不占学院普通名额。
  \item 接收端要求：国防科技大学《接收2027级地方院校推荐免试硕士研究生（含直博生、参军入伍生）工作方案》要求申请人「取得所在学校硕士研究生推荐免试计划（\textbf{含我校『国防科工计划』}）」；另设「参军入伍生」类别，须符合军队直接选拔招录军官条件（本科院校为「双一流」建设高校、指定学科门类等），\textbf{具体政策以当年军队通知为准}。
  \item 政审：该校要求申请人「具有良好的思想品德和政治素质……\textbf{有为国防和军队现代化建设服务的思想基础}」；军事院校另按军队规定组织政治考核。
  \item 待遇佐证：东华理工大学2026年招生简章写明「国防科工院校补偿计划硕士研究生享受专项奖学金，奖励标准为8000元/生·学年」。
\end{itemize}
\nref{S0561}
\nref{S0650}
\nref{S0420}
\nref{S0602}
\end{qa}

\iin{南开大学}
\ith{军工定向}
\ist{特殊通道与专项}
\begin{qa}{E0236}{定向生违约要赔多少钱？怎么算？}
结论：\textbf{本所查公开资料中没有检索到全国统一的「定向生违约金计算标准」；能确认的是违约要退还已享受的定向培养费用并缴纳违约金，数额在个人协议里约定。}
\begin{itemize}
  \item 有明确文件依据的是公费师范生（国办发〔2024〕27号）：第十四条「\textbf{按违约处理，须退还已享受的公费教育费用并缴纳违约金}。教育部、财政部结合学生违约退费等情况，确定学校年度补助资金安排」；第十三条「未履约任教或履约任教期限不足，按协议约定承担相应责任，\textbf{违约行为纳入信用记录}，并及时与全国信用信息共享平台共享，\textbf{违约相关材料归入本人人事档案}……完成履约任务后即修复相关信用」。
  \item 优师计划同理（教师〔2021〕4号）：未按规定履约的「按照协议约定承担相应责任，\textbf{违约记录归入人事档案，依法依规纳入社会信用体系}」。
  \item 农村订单定向医学生：服务期内报考受限（电子科技大学2026年章程规定只能报考临床医学专业学位硕士全科医学方向，须提供定向县卫健部门书面同意函，违约后果由考生本人承担）。
\end{itemize}
\textbf{签约前请逐条读清协议条款，以上为文件规定，个人协议可能另有约定。}
\nref{S0456}
\nref{S0543}
\nref{S0493}
\end{qa}

\iin{教育部学生服务与素质发展中心（学信网/研招网）}
\ith{定向违约}
\ist{特殊通道与专项}
\section{支教团}
\begin{qa}{E0213}{研究生支教团是什么？怎么报名？服务一年后直接读研吗？}
研究生支教团=\textbf{先取得推免资格、保留研究生入学资格，到中西部中小学支教一年，服务期满考核合格后返校读研}，属于推免的一种，不是「读完再考」。
\begin{itemize}
  \item 主体与时长：面向「具备推免资格的应届本科毕业生和\textbf{在读研究生}」，到对口帮扶地区中小学开展一年支教（2027年8月—2028年7月）（安徽工业大学第29届方案、西南交通大学2027年通知）；西南交大2026年招募18人，服务地为甘肃秦安、四川马尔康、新疆莎车等。
  \item 报名流程（安徽工业大学方案）：学生报名→学院推荐（每学院不多于2名）→学校选拔（学业成绩20\%＋面试80\%，差额面试、全程录像、纪检监督）→体检→公示不少于7日→\textbf{9月30日前}在学生推免系统完成信息录入并签《招募协议书》。
  \item 依据：共青团中央办公厅、教育部办公厅2026年4月21日印发的《关于进一步加强中国青年志愿者研究生支教团工作的意见》及《中国青年志愿者研究生支教团管理细则》（全国项目办发〔2026〕2号）。
\end{itemize}
\nref{S0589}
\nref{S0061}
\nref{S0370}
\nref{S0614}
\end{qa}

\iin{安徽工业大学}
\ith{支教团}
\ist{特殊通道与专项}
\begin{qa}{E0214}{支教团的名额和普通推免名额是分开的吗？算不算推免生？}
结论：\textbf{名额分开、不占用学校普通推免指标；支教团成员本身就是推免生。}
\begin{itemize}
  \item 指标来源（安徽工业大学2027年方案）：「第29届招募指标由\textbf{教育部专项下拨}为4名。研究生支教团推免指标\textbf{不占用学校计划推荐免试研究生的名额}。」西南交通大学通知同：「该推免指标由教育部专项下拨，不占用学校下达给各学院的研究生免试推荐名额」；郑州大学通知亦同。
  \item 身份认定（华中师范大学2027年通知）：支教团招募是「\textbf{推荐优秀应届本科毕业生免试攻读硕士研究生工作的重要组成部分}」，被推荐学生的学业综合成绩基本要求应与所在学科专业普通类推免生保持一致。
  \item 名目归属：纳入西部计划\textbf{乡村教育专项}统筹实施；在研招网体系对应独立的「\textbf{支教团推免计划}」校验码；湖南大学亦将「研究生支教团等\textbf{国家专项计划}推免生名额按国家相应政策分配至相关推荐单位」。
\end{itemize}
\nref{S0589}
\nref{S0061}
\nref{S0597}
\nref{S0614}
\nref{S0540}
\end{qa}

\iin{安徽工业大学}
\ith{支教团}
\ist{特殊通道与专项}
\begin{qa}{E0215}{支教团选拔什么时候开始？要笔试面试吗？}
结论：\textbf{选拔集中在每年8月底至9月底（与推免同步），以面试为主，要求差额面试、全程录像并接受纪检监督。}
\begin{itemize}
  \item 时间表（安徽工业大学第29届方案）：9月11日16:00前学院报送推荐材料→学校资格复审与遴选→统一体检→公示不少于7日→\textbf{9月30日前}在学生推免系统完成信息录入并签《招募协议书》。
  \item 考核构成：总成绩=学业成绩20\%＋面试80\%。学业成绩按专业排名打分（前5\%为100分、5\%—20\%为95分、20\%—35\%为90分）；面试为差额面试，自我陈述与答题各占50\%，去掉最高分和最低分后取平均。
  \item 报名门槛：学业成绩\textbf{专业排名前35\%}，完成前三年（五年制前四年）培养计划课程（不含公共选修课）及第二课堂积分，四级达到国家规定报考六级分数要求（艺术设计类、高水平运动员不作硬性要求），有教师资格证者同等条件下优先。
  \item 纪律要求（共青团中央、教育部《意见》）：集中面试等选拔过程应\textbf{全程录音录像}，资料长期留存备查。
\end{itemize}
\nref{S0589}
\nref{S0061}
\nref{S0370}
\nref{S0614}
\end{qa}

\iin{安徽工业大学}
\ith{支教团}
\ist{特殊通道与专项}
\begin{qa}{E0216}{支教团服务期间算不算应届？服务期满后还能考公吗？}
结论：\textbf{服务期间不算应届}（你已本科毕业，只是保留研究生入学资格）；服务期满考核合格可享受西部计划服务期为一年的政策，含计算服务期与工龄。\textbf{「考公是否算应届」本所查公开资料没有教育部统一规定。}
\begin{itemize}
  \item 在校状态（安徽工业大学2027年方案）：「服务期间，推免录取的学生\textbf{保留研究生入学资格}，在读研究生\textbf{保留学籍}」，人事关系挂靠学校，党团组织关系临时转到服务地（郑州大学通知）。
  \item 可确认待遇：安徽工业大学方案明确「服务期满且\textbf{考核合格}的志愿者，享受西部计划服务期为\textbf{一年}的有关政策，\textbf{依实际服务年限计算服务期及工龄}，并在服务证书及服务鉴定表中体现」；郑州大学通知写「一年的支教服务工作时间在日后工作的\textbf{工龄计算}中有效……考核合格者由团中央等部门颁发《志愿者证书》，此证书可作为\textbf{工龄证明}依据」。
  \item 保障：西部计划政策问答明确服务省项目办为志愿者（\textbf{含研究生支教团志愿者}）购买重大疾病、人身意外伤害等商业保险；西部计划服务期为1至3年、协议一年一签。
  \item 关于考公：教育部、团中央文件中\textbf{未见}「支教团服务期满按应届生对待」的统一表述，请以当年目标省招录公告的界定为准。
\end{itemize}
\nref{S0589}
\nref{S0597}
\nref{S0598}
\nref{S0370}
\end{qa}

\iin{安徽工业大学}
\ith{支教团}
\ist{特殊通道与专项}
\begin{qa}{E0217}{支教团服务期考核不合格会怎样？还能不能读研？}
结论：\textbf{考核不合格会被取消研究生入学资格；录取后中途退出的，还会连带影响所在学院下一届推荐指标。}
\begin{itemize}
  \item 考核不合格（郑州大学第二十四届通知）：「考核工作由国家、省级西部计划项目办和我校支教团工作领导小组共同负责。学生返校入学后一个月内公布考核结果……分优秀、合格、不合格三类。\textbf{考核不合格者，取消其研究生入学资格。}」
  \item 录取后主动退出（华中师范大学2027年通知，援引《中国青年志愿者研究生支教团管理细则》全国项目办发〔2026〕2号）：志愿者录取后因升学、就业、出国、个人健康等主客观原因退出，「领导小组办公室将按相关规定\textbf{取消其研究生支教团志愿者资格，并酌情核减学生所在学部学院下一年度推荐指标或取消下一年度推荐资格}」。
  \item 服务期保障：按《大学生志愿服务西部计划志愿者管理办法》与《管理细则》执行，工作生活补贴、社会保险、商业保险按当年度西部计划实施方案落实（安徽工业大学2027年方案）。
  \item 推荐标准不降：共青团中央、教育部《意见》要求支教团推免生学业综合成绩基本要求与所在学科专业普通类推免生保持一致。
\end{itemize}
\nref{S0597}
\nref{S0614}
\nref{S0370}
\nref{S0589}
\end{qa}

\iin{郑州大学}
\ith{支教团}
\ist{特殊通道与专项}
\begin{qa}{E0264}{卓越工程师学院、工程硕博士培养改革专项是什么？怎么报？}
结论：\textbf{它是国家层面推动的「工程硕博士培养改革专项」，由高校卓越工程师学院/相关学院承办，采取校企联合培养、工学交替模式，通过推免渠道招收硕士和直博。}
\begin{itemize}
  \item 定位与模式（中国人民大学2027年专项推免研究生招生报名通知）：「以培养卓越工程师后备人才为目标，聚焦国家重大战略需求，支撑产业链安全……专项分为\textbf{工程硕士和工程直博两个项目，实行联合培养和工学交替模式，校企导师组共同承担研究生培养工作}。」招生领域为电子信息（085400）人工智能方向，合作企业含腾讯、华为、百度、智谱华章、小米、北汽等，各学院各方向计划数逐条列出。
  \item 学制（同通知）：硕士基本学制3年（约1年校内课程＋2年联合培养单位实践），直博5年（约2年校内＋3年实践）；录取者均为全日制\textbf{非定向}工程类专业学位研究生。
  \item 报名（同通知）：登录学校推免生信息采集系统填报，每位申请者\textbf{限报一个专业方向，不能同时申请硕士项目和直博项目}；学院考核选拔后，\textbf{最终仍须按教育部规定时间在推免服务系统完成填报志愿、复试确认与待录取确认}。
  \item 常见并列专项（东南大学2027年推免专项报名通知）：工程硕博士、集成电路专项班、国优计划、研支团四个专项，「报名时只能选择一个专项，不可兼报」。
\end{itemize}
\nref{S0286}
\nref{S0512}
\nref{S0188}
\nref{S0186}
\nref{S0548}
\end{qa}

\iin{中国人民大学}
\ith{工程专项}
\ist{特殊通道与专项}
\begin{qa}{E0265}{工程硕博士专项和普通推免怎么选？哪个更值？}
结论：\textbf{没有绝对更值，取决于你要学术路径还是工程/产业路径；制度上它是专项、名额相对独立，且与普通推免互斥。}
\begin{itemize}
  \item 互斥性（东南大学2027年推免专项报名通知）：四个专项「报名时\textbf{只能选择一个专项，不可兼报}……提交后不能再做更改」；「\textbf{确认拟录取的考生，不得再参选其他任何专项}」。
  \item 培养差异（中国人民大学2027年专项通知）：实行「联合培养和工学交替模式，校企导师组共同承担研究生培养工作，学生深入工程攻关一线」；硕士学制3年（1年校内课程＋2年左右企业实践），直博5年（2年左右校内＋3年左右实践）；录取者均为全日制\textbf{非定向工程类专业学位}研究生——学位类型是专业学位，不是学术学位。
  \item 名额性质：南开大学推免实施办法把「定向补偿推免（支教团、硕师计划、国防科工院校补偿计划、新疆高层次双语人才计划等）」单列、由牵头部门实施；湖南大学推免通知规定「研究生支教团等\textbf{国家专项计划}推免生名额按国家相应政策分配至相关推荐单位」，即专项名额通常不占学院普通名额。
  \item 普通推免的优势：学位类型、导师、方向自由度更高，可申学术型硕士或直博，也可内推外推。
\end{itemize}
\nref{S0512}
\nref{S0286}
\nref{S0548}
\nref{S0561}
\nref{S0540}
\end{qa}

\iin{东南大学}
\ith{工程专项}
\ist{特殊通道与专项}
\section{公费师范}
\begin{qa}{E0229}{农村专项或者国家专项计划录取的本科生，能正常保研吗？}
结论：\textbf{能。国家专项、地方专项、高校专项都属于普通高考招生录取，不构成推免资格障碍；被排除的是专升本、第二学士学位和定向就业生。}
\begin{itemize}
  \item 资格口径（多校一致的模板表述）：推荐对象为「本校通过\textbf{普通高考（含保送方式）招生录取}或面向港澳台华侨学生招生录取的全日制应届本科毕业生，\textbf{不含专升本、第二学士学位学生}」——燕山大学2027年通知、北京科技大学2027年通知、辽宁大学2027年办法、成都大学商学院2027年细则均如此。专项计划录取依据仍是普通高考成绩与招生计划。
  \item 名额测算上的排除项是身份类别而非批次：安徽师范大学2027年方案的基础指标按「全校2027届本科预计毕业学生（\textbf{不含专升本与定向就业生}）」总数折算；浙江师范大学2027年通知的推荐对象「不含定向生，但原定向单位书面同意推荐或硕师计划有关教育局指定的除外」。
  \item 需另作区分：若入学时签过定向就业协议（如农村订单定向医学生），那是\textbf{定向生}身份问题——中山大学2027年办法明确不接收未按国家规定完成履约责任（含违约）的「农村订单定向生」。
\end{itemize}
\nref{S0533}
\nref{S0233}
\nref{S0646}
\nref{S0569}
\nref{S0166}
\end{qa}

\iin{zsjyc（高校官网）}
\ith{农村专项}
\ist{特殊通道与专项}
\begin{qa}{E0239}{公费师范生“本硕衔接”或者 4+2 是怎么选拔的？}
结论：\textbf{选拔方式=学校对本科前三年学业成绩与在校表现做「推免综合考核」，取得推免资格并达到硕士录取标准的可推免录取本校全日制教育硕士；所需推免名额统筹下达至学校总推免名额中。}
\begin{itemize}
  \item 制度原文（国办发〔2024〕27号第七条）：「学校结合\textbf{本科前三年学业成绩、在校表现}等对学生进行\textbf{推荐免试攻读研究生综合考核}，取得学校推免资格并达到学校硕士研究生录取标准的可推免录取至本校全日制教育硕士专业学位研究生，并于\textbf{本科第四年开始硕士课程先修学习}，所需推免名额\textbf{统筹下达至学校总推免名额中}。退出公费师范生培养的学生，不再具有申请推免资格。」
  \item 培养结构（第八条）：「公费师范生\textbf{本科4年及教育硕士研究生2年}培养方案，实行\textbf{分流淘汰}」——这就是「4+2」的来源。
  \item 过渡安排（教师厅函〔2025〕3号）：该办法印发前已签协议的部属师大本科一至三年级在读者（2025—2027届），\textbf{本科三年级结束后}可自愿申请，经学校推免综合考核并达到录取标准，\textbf{报生源省份省级教育行政部门同意并签署补充协议}后免试攻读本校全日制教育硕士再履约任教。
  \item 分流节点（第十条）：\textbf{研究生一年级课程学习结束后}，学校根据本科以来综合考核结果排序，据此在公布范围内选择履约任教地（市、州、盟）。
\end{itemize}
\nref{S0456}
\nref{S0044}
\nref{S0559}
\end{qa}

\iin{教育部学生服务与素质发展中心（学信网/研招网）}
\ith{公费师范}
\ist{特殊通道与专项}
\begin{qa}{E0240}{优师计划能不能读研？服务期内读研算违约吗？}
结论：\textbf{优师计划本身是本科层次定向培养；2027年起部分省份已把「优师计划」与「硕师计划」结合实施，优师计划本科应届毕业生可经此渠道推免攻读教育硕士。服务期内能否再读研，须按协议约定。}
\begin{itemize}
  \item 原始定位（教师〔2021〕4号，教育部等九部门《中西部欠发达地区优秀教师定向培养计划》）：须参加当年全国统一高考、\textbf{本科提前批次}录取；「国家优师专项师范生在录取后、获得录取通知书前，须与培养学校和生源所在省份省级教育行政部门、乡村振兴工作部门\textbf{签订协议，承诺毕业后到生源所在省份定向县中小学履约任教不少于6年}。对拒签协议的录取学生，取消其优师专项录取资格。」
  \item 新增衔接通道（吉林师范大学2027年硕师计划通知）：「从2027年起，省级主管部门将『\textbf{优师计划』与『硕师计划』结合实施}」；该校2027年获推免指标80个（常规硕师计划76个、优师计划专项4个），「『优师计划』专项推荐对象为……『\textbf{优师计划』本科应届毕业生}」。
  \item 违约后果（教师〔2021〕4号）：未按规定履约的「按照协议约定承担相应责任，\textbf{违约记录归入人事档案，依法依规纳入社会信用体系}」。
  \item 本所查公开资料\textbf{未检索到}「优师计划服务期内读研」的统一条款。
\end{itemize}
\textbf{签约前请逐条读清协议条款，以下为文件规定，个人协议可能另有约定。}
\nref{S0543}
\nref{S0593}
\nref{S0592}
\end{qa}

\iin{中华人民共和国教育部}
\ith{公费师范}
\ist{特殊通道与专项}
\begin{qa}{E0241}{公费师范生违约的代价是什么？会进征信吗？}
结论：\textbf{会。违约要退还已享受的公费教育费用并缴纳违约金，违约行为纳入信用记录并与全国信用信息共享平台共享，违约材料归入人事档案。}
\begin{itemize}
  \item 制度原文（国办发〔2024〕27号第十三条）：「公费师范生毕业后均须履约任教，未履约任教或履约任教期限不足，按《本研衔接师范生公费教育协议》约定承担相应责任，\textbf{违约行为纳入信用记录}，并及时与全国信用信息共享平台共享，\textbf{违约相关材料归入本人人事档案}。履约任务完成前\textbf{不得到定向地（市、州、盟）之外从教，不得被招录（聘用）为其他公务人员}。\textbf{完成履约任务后即修复相关信用。}」
  \item 违约触发情形（第十四条）：除生病、应征入伍等可申请中止协议的情形，以及「因学业水平未达到免试攻读教育硕士研究生条件」等可按协议解除的情形外，「未在学校规定学习年限内取得教育硕士研究生学历学位证书和教师资格证书的，\textbf{按违约处理，须退还已享受的公费教育费用并缴纳违约金}」。
  \item 同类机制：优师计划「未按规定履约的，按照协议约定承担相应责任，违约记录归入人事档案，依法依规纳入社会信用体系」（教师〔2021〕4号）。
\end{itemize}
\textbf{签约前请逐条读清协议条款，以上为文件规定，个人协议可能另有约定。}
\nref{S0456}
\nref{S0044}
\nref{S0543}
\end{qa}

\iin{教育部学生服务与素质发展中心（学信网/研招网）}
\ith{公费师范}
\ist{特殊通道与专项}
\begin{qa}{E0267}{部属师范大学的“本研衔接”公费师范生政策怎么理解？}
结论：\textbf{核心是「本科4年＋教育硕士2年」一贯培养、免试攻读本校全日制教育硕士、毕业后回生源省定向地市任教不少于6年，实行分流淘汰，退出者不再具有推免资格。}
\begin{itemize}
  \item 文件依据（国办发〔2024〕27号《教育部直属师范大学本研衔接师范生公费教育实施办法》，2024年5月28日国务院办公厅转发，同时废止2018年公费师范教育实施办法）：第七条明确培养目标为全日制教育硕士专业学位研究生；第八条明确「统筹设计并实施公费师范生\textbf{本科4年及教育硕士研究生2年}培养方案，实行\textbf{分流淘汰}」。
  \item 履约要求（第九条）：三方签订《本研衔接师范生公费教育协议》，「公费师范生毕业后一般回生源所在省份\textbf{定向地（市、州、盟）中小学任教，并承诺从事中小学教育工作6年以上}。到城镇学校工作的，应到农村义务教育学校任教服务\textbf{至少1年}」。
  \item 定向地确定（第十条）：\textbf{研一课程学习结束后}，学校按本科以来综合考核结果排序，公费师范生按排序在公布范围内选择任教地。
  \item 退出机制（第七条）：\textbf{退出公费师范生培养的学生，不再具有申请推免资格}；入学1年内经考察不适合从教的，退还已享受的公费教育费用并调整到非师范专业。
  \item 过渡安排（教师厅函〔2025〕3号）：2025—2027届本科层次公费师范生在本科三年级结束后可自愿申请转入本研衔接，经推免综合考核、报省级教育行政部门同意并签署补充协议后生效。
\end{itemize}
\nref{S0456}
\nref{S0044}
\nref{S0559}
\end{qa}

\iin{教育部学生服务与素质发展中心（学信网/研招网）}
\ith{公费师范}
\ist{特殊通道与专项}
\section{退役免试}
\begin{qa}{E0218}{退役大学生士兵免初试要满足什么条件？}
结论：\textbf{条件是：服役期间获得三等战功、二等功以上奖励或二级以上表彰，且符合全国硕士研究生招生考试报考条件。}
\begin{itemize}
  \item 法规原文（《2027年全国硕士研究生招生工作管理规定》第二十五条；2026年第二十五条同）：「对服役期间获得\textbf{三等战功、二等功以上奖励或者二级以上表彰}，符合全国硕士研究生招生考试报考条件的退役人员，可申请免初试攻读硕士研究生。符合免初试资格的考生，应在教育部规定的\textbf{全国统考报名时间内}登录『全国推荐免试攻读研究生（免初试、转段）信息公开暨管理服务系统』报名，\textbf{逾期不得补报}。招生单位对审核通过的考生组织复试，\textbf{复试合格者即可予以拟录取}。」
  \item 审核（第二十六条）：须在报名时填报申请信息，有关部门、报考点初审，招生单位\textbf{复试前复审}；未按规定申报或审核未通过的，\textbf{不享受}免初试政策。
  \item 再报统考的限制（推免服务系统说明第（五）条）：「申请立功表彰免初试的考生，\textbf{不得再报名参加全国硕士研究生招生考试}。」
  \item 材料：立功受奖证明，以及《入伍批准书》《退出现役证》（限士兵）等。
\end{itemize}
\nref{S0442}
\nref{S0437}
\nref{S0051}
\nref{S0464}
\nref{S0607}
\end{qa}

\iin{教育部学生服务与素质发展中心（学信网/研招网）}
\ith{退役免试}
\ist{特殊通道与专项}
\begin{qa}{E0219}{三等功、二等功免初试怎么申请？系统里走哪个入口？}
结论：\textbf{入口是「全国推荐免试攻读研究生（免初试、转段）信息公开管理服务系统」（https://yz.chsi.com.cn/tm），须在统考报名期内报名，逾期不得补报。}
\begin{itemize}
  \item 系统与时间（《2027年全国硕士研究生招生工作管理规定》第二十五条）：「符合免初试资格的考生，应在教育部规定的\textbf{全国统考报名时间内}登录『推免服务系统』报名，\textbf{逾期不得补报}。招生单位对审核通过的考生组织\textbf{复试}，复试合格者即可予以拟录取。」2027年统考网上报名为2026年10月15日—10月24日（预报名10月9日—12日）。
  \item 系统说明对应条款：「服役期间获得三等战功、二等功以上奖励或者二级以上表彰，符合全国硕士研究生招生考试报考条件的退役人员，可申请免初试……\textbf{申请立功表彰免初试的考生，不得再报名参加全国硕士研究生招生考试。}」
  \item 资格审核（第二十六条）：报名时填报申请信息，有关部门初审、招生单位在\textbf{复试前复审}；未按规定申报或审核未通过的，不享受免初试政策。
  \item 材料：立功受奖证明、退役证件等。
\end{itemize}
\nref{S0442}
\nref{S0437}
\nref{S0464}
\nref{S0607}
\end{qa}

\iin{教育部学生服务与素质发展中心（学信网/研招网）}
\ith{退役免试}
\ist{特殊通道与专项}
\begin{qa}{E0220}{退役士兵专项计划和立功免初试是一回事吗？}
结论：\textbf{不是一回事。前者是统考招生计划（要参加初试、招生单位自主划线），后者是免初试（不考初试、直接复试）。}
\begin{itemize}
  \item 退役大学生士兵专项计划（教学厅〔2025〕8号）：「在硕士研究生招生计划总规模内\textbf{单列下达、专项使用}，不得调整为普通计划或其他专项计划」；2026年由北京大学、清华大学等\textbf{532所}高校承担；招生单位\textbf{自主确定}该计划考生初试成绩要求；报考者须为「高校学生应征入伍退出现役且符合全国硕士研究生报考条件者」。
  \item 立功免初试（2027年管理规定第二十五条）：面向服役期间获\textbf{三等战功、二等功以上奖励或二级以上表彰}的退役人员，\textbf{免初试}、由招生单位组织复试，复试合格即可拟录取，且报名后不得再参加统考。
  \item 系统位置不同：专项计划在\textbf{全国统考}考试方式下填报（第二十四条列举「退役大学生士兵计划」）；免初试在\textbf{推免服务系统}的「免初试」入口。
  \item 互斥限制（第二十四、二十六条）：报考专项计划的考生「\textbf{不再享受初试加分、少数民族照顾政策}」；填报后资格复核未通过的，不得按该计划复试，仅可在调剂阶段申请调剂到普通计划。
\end{itemize}
\nref{S0052}
\nref{S0442}
\nref{S0437}
\nref{S0440}
\end{qa}

\iin{中华人民共和国教育部}
\ith{计划辨析}
\ist{特殊通道与专项}

\begin{qa}{N0006}{强基计划本研转段为什么也出现在推免服务系统里？}
系统名称中的“免初试、转段”同时覆盖普通推免和部分转段信息备案，不代表两类学生走完全相同的招生流程。研招网2027届说明单列：2023级拟转段强基生于2026年10月13日至20日查询并确认本人转段信息，确认结果作为进入研究生阶段的录取信息。强基生应先完成本校培养方案和转段考核，再按系统提示核对；不要误填普通推免志愿，也不要把该届日期套用到其他年级。\nref{S0761}
\qscope{全国推免服务系统2027届说明；该日期仅适用于说明所列2023级拟转段强基生。}
\end{qa}

\begin{qa}{N0007}{立功表彰退役人员申请免初试，和普通推免生走的是同一条报名流程吗？}
两者都使用推免服务系统，但资格条件、入口和报名窗口并不相同。2027年招生管理规定规定，符合条件的退役人员可申请免初试攻读硕士，需在全国统考报名时间内通过该系统报名，逾期不得补报；系统说明还提示申请此项免初试者不得再报名参加当年全国统考。报名之前应按规定准备立功表彰和退役材料，并向招生单位确认复试安排。\nref{S0762}
\qscope{教育部《2027年全国硕士研究生招生工作管理规定》第二十五条（2026年发布）；推免服务系统2027届说明。}
\end{qa}
\begin{qa}{E0259}{立功免试是不是必须在统考报名期内报名？过期能补报吗？}
结论：\textbf{必须在统考报名期内报名，逾期不得补报——这是文件里写死的硬截止。}
\begin{itemize}
  \item 文件原文（《2027年全国硕士研究生招生工作管理规定》第二十五条）：「符合免初试资格的考生，应在教育部规定的\textbf{全国统考报名时间内}登录『全国推荐免试攻读研究生（免初试、转段）信息公开暨管理服务系统』（网址：https://yz.chsi.com.cn/tm）报名，\textbf{逾期不得补报}。」《2026年管理规定》第二十五条为同一表述。
  \item 时间锚点（2027年管理规定第十九条）：2027年统考网上报名为\textbf{2026年10月15日至10月24日}，预报名为10月9日至10月12日，每日9:00—22:00；且「所有考生须按要求如实填写报名信息并提供真实材料……按时提交补充材料，\textbf{逾期不得补办}」。
  \item 系统侧提示（推免服务系统说明第（五）条）：「符合免初试资格的考生，应在教育部规定的全国统考报名时间内登录『推免服务系统』报名，\textbf{逾期不得补报}。申请立功表彰免初试的考生，\textbf{不得再报名参加全国硕士研究生招生考试}。」
  \item 后续环节：招生单位对审核通过的考生组织复试，复试合格者即可予以拟录取（第二十五条）；资格由有关部门初审、招生单位复试前复审（第二十六条）。
\end{itemize}
\nref{S0442}
\nref{S0437}
\nref{S0464}
\nref{S0051}
\end{qa}

\iin{教育部学生服务与素质发展中心（学信网/研招网）}
\ith{退役免试}
\ist{特殊通道与专项}
